\documentclass{amsart}
\usepackage{amsmath,amssymb,amsthm,hyperref}
\newtheorem{theorem}{Theorem}
\newtheorem{lemma}{Lemma}
\begin{document}

\begin{center}\large\textbf{Optimal detector configuration for $S/\sqrt{B}$}\end{center}

\section*{Setup}

Index the three independent modifications by $i\in\{1,2,3\}$:
\begin{itemize}
\item $i=1$: reduced boost $\beta\gamma=0.28$ (vs.\ the baseline $0.56$);
\item $i=2$: inclusion of the forward PID device (Fwd-PID);
\item $i=3$: inclusion of the backward EMC (Bwd-EMC).
\end{itemize}
A \emph{configuration} is a vector $x=(x_1,x_2,x_3)\in\{0,1\}^3$, where $x_i=1$
means modification $i$ is active. There are $2^3=8$ configurations.

Each modification, when active, multiplies the signal yield $S$ and the
background yield $B$ by fixed positive factors. Writing $(s_i,b_i)$ for the
(signal, background) multipliers of modification $i$, the empirical data are
\[
(s_1,b_1)=(1.07,\;0.94),\qquad
(s_2,b_2)=(1.025,\;1.025),\qquad
(s_3,b_3)=(1.00,\;0.90).
\]
The baseline configuration is $x=(0,0,0)$; by definition its signal and
background are $S_0>0$ and $B_0>0$.

\paragraph{Independence assumption.}
The modifications act independently and multiplicatively, so for a
configuration $x$,
\[
S(x)=S_0\prod_{i=1}^{3}s_i^{\,x_i},
\qquad
B(x)=B_0\prod_{i=1}^{3}b_i^{\,x_i}.
\]
(With $x_i\in\{0,1\}$, the factor for modification $i$ is $s_i$ if $x_i=1$ and
$1$ if $x_i=0$, i.e.\ $s_i^{x_i}$; likewise for $b_i$.) Since all $s_i,b_i>0$
and $S_0,B_0>0$, we have $S(x)>0$, $B(x)>0$ for every $x$.

The figure of merit is
\[
F(x)\;=\;\frac{S(x)}{\sqrt{B(x)}}.
\]
We seek $\arg\max_{x\in\{0,1\}^3}F(x)$ and the improvement factor
$R(x)=F(x)/F(0,0,0)$.

\section*{The improvement factor decomposes into independent per-option gains}

\begin{lemma}\label{lem:decomp}
For every $x\in\{0,1\}^3$,
\[
R(x)\;=\;\frac{F(x)}{F(0,0,0)}\;=\;\prod_{i=1}^{3}g_i^{\,x_i},
\qquad\text{where}\qquad
g_i:=\frac{s_i}{\sqrt{b_i}}>0 .
\]
\end{lemma}

\begin{proof}
Using the multiplicative forms of $S(x)$ and $B(x)$ and the fact that
$\sqrt{\cdot}$ is multiplicative on positive reals,
\[
F(x)=\frac{S_0\prod_i s_i^{x_i}}{\sqrt{B_0\prod_i b_i^{x_i}}}
=\frac{S_0}{\sqrt{B_0}}\;\prod_{i=1}^{3}\frac{s_i^{x_i}}{\sqrt{b_i^{x_i}}}
=F(0,0,0)\,\prod_{i=1}^{3}\Bigl(\frac{s_i}{\sqrt{b_i}}\Bigr)^{x_i}.
\]
(Here $b_i^{x_i}\ge 0$ and $\sqrt{b_i^{x_i}}=(\sqrt{b_i})^{x_i}$ for
$x_i\in\{0,1\}$.) Dividing by $F(0,0,0)>0$ gives the claim with
$g_i=s_i/\sqrt{b_i}$.
\end{proof}

\paragraph{Per-option net gains.}
Numerically,
\[
g_1=\frac{1.07}{\sqrt{0.94}}=1.103620\ldots,\quad
g_2=\frac{1.025}{\sqrt{1.025}}=\sqrt{1.025}=1.012422\ldots,\quad
g_3=\frac{1.00}{\sqrt{0.90}}=1.054092\ldots .
\]
In particular each $g_i>1$. We verify this exactly:
\begin{itemize}
\item $g_1>1 \iff 1.07>\sqrt{0.94} \iff 1.07^2=1.1449>0.94$. True.
\item $g_2=\sqrt{1.025}>1$ since $1.025>1$. True.
\item $g_3>1 \iff 1>\sqrt{0.90} \iff 1>0.90$. True.
\end{itemize}

\section*{Optimization}

\begin{theorem}\label{thm:main}
The figure of merit $F$ over the eight configurations is uniquely maximized by
the all-on configuration $x^\star=(1,1,1)$ (reduced boost, Fwd-PID, and
Bwd-EMC all active). The maximal improvement factor relative to the baseline is
\[
R(x^\star)=g_1g_2g_3
=\frac{1.07\cdot 1.025}{\sqrt{0.94\cdot 1.025\cdot 0.90}}
=\frac{1.09675}{\sqrt{0.867150}}\;=\;1.17777\ldots,
\]
i.e.\ about a $17.8\%$ improvement.
\end{theorem}

\begin{proof}
By Lemma~\ref{lem:decomp}, maximizing $F(x)$ over $x\in\{0,1\}^3$ is
equivalent to maximizing $R(x)=\prod_{i}g_i^{x_i}$, because
$F(x)=F(0,0,0)\,R(x)$ and $F(0,0,0)>0$ is a constant independent of $x$.

We show $x^\star=(1,1,1)$ is the strict global maximizer directly. Fix any
$x\neq x^\star$. Then the set $Z:=\{\,j:\ x_j=0\,\}$ of inactive options is
nonempty. Using Lemma~\ref{lem:decomp},
\[
\frac{R(x^\star)}{R(x)}
=\frac{\prod_{i=1}^3 g_i^{\,1}}{\prod_{i=1}^3 g_i^{\,x_i}}
=\prod_{i=1}^3 g_i^{\,1-x_i}
=\prod_{j\in Z} g_j .
\]
Each factor satisfies $g_j>1$ (shown above) and $Z\neq\varnothing$, so the
product exceeds $1$; hence $R(x^\star)>R(x)$, equivalently $F(x^\star)>F(x)$.
As $x\neq x^\star$ was arbitrary, $x^\star$ is the unique maximizer of $F$ over
all eight configurations.

Finally,
\[
R(x^\star)=g_1g_2g_3=\frac{s_1 s_2 s_3}{\sqrt{b_1 b_2 b_3}}
=\frac{1.07\cdot 1.025\cdot 1.00}{\sqrt{0.94\cdot 1.025\cdot 0.90}}.
\]
The numerator is $1.07\cdot1.025=1.09675$. The radicand is
$0.94\cdot1.025\cdot0.90=0.867150$, with $\sqrt{0.867150}=0.931209\ldots$, so
\[
R(x^\star)=\frac{1.09675}{0.931209\ldots}=1.177770\ldots,
\]
as claimed.
\end{proof}

\section*{Remark on robustness}

The structural conclusion is independent of the precise numerical values: an
option should be switched on \emph{iff} its net gain $g_i=s_i/\sqrt{b_i}$
exceeds $1$, equivalently iff $s_i^2>b_i$ (signal-squared improvement beats
background inflation). Here all three options satisfy $s_i^2>b_i$
($1.1449>0.94$, $1.050625>1.025$, $1>0.90$), which is why the all-on
configuration is optimal. The reduced boost contributes the largest single
gain ($\approx 10.4\%$), the Bwd-EMC next ($\approx 5.4\%$), and the Fwd-PID
the smallest ($\approx 1.2\%$): note the Fwd-PID multiplies $S$ and $B$ by the
\emph{same} factor $1.025$, yet still helps because $S/\sqrt B$ scales as
$1.025/\sqrt{1.025}=\sqrt{1.025}>1$.

\medskip
\noindent\textbf{Answer.} The optimal configuration is
\emph{reduced boost $+$ Fwd-PID $+$ Bwd-EMC} (all three on), giving an
$S/\sqrt B$ improvement factor of
$\dfrac{1.07\cdot1.025}{\sqrt{0.94\cdot1.025\cdot0.90}}\approx 1.178$
(about $+17.8\%$) over the baseline.

\end{document}
